% ==============================================================================
% ПАРТИЯ B03-026: DISS-005-B1470 .. DISS-005-B1524
% Глава 3: Введение, методология формирования датасета и вызовы
% Замечания: DISS-005-C0103 .. DISS-005-C0111 (9 шт.)
% Объекты: 0 шт.
% ==============================================================================

% DISS-005-B1471
\section{Введение}\label{sec:ch3_intro}

% DISS-005-B1472
Технологические\commentnote{DISS-005-C0104}{Алексей Евдаков [2]}{Взято из черновика статьи «Development\_of\_an\_Open\_Dataset\_of\_Real\_World\_Oscillograms\_from\_Electric\_Power\_Grids»} достижения последних столетий стимулировали создание и развитие электроэнергетических систем, обеспечивающих стабильное электроснабжение различных секторов экономики и поддерживающих высокий уровень жизни населения [51, 52, 53, 54, 55]. На заре электроэнергетики прямая запись осциллограмм электрических сигналов была невозможна. Вместо этого анализ электрических явлений основывался на теоретических моделях и экспертных оценках [56]. Глубокие исследования многочисленных процессов требовали обширных экспериментов и наблюдений, что ограничивало возможности анализа и предотвращения аварийных ситуаций.\par\vspace{0.2cm}

% DISS-005-B1473
По мере расширения и усложнения энергосистем все более актуальными становились вопросы качества электроэнергии и стабильности электроснабжения [57, 58]. Рост потребления электроэнергии, разнообразие потребителей и интеграция новых промышленных технологий сделали систему более уязвимой к различным сбоям и нарушениям [59, 60]. Для решения этих проблем были усовершенствованы устройства релейной защиты и автоматики (РЗА) [56], ключевой функцией которых является защита и управление основным оборудованием энергосистемы путем реагирования на превышение контролируемыми параметрами заданных порогов (уставок). Ранние устройства РЗА использовали электромеханические реле [61], которые под действием электромагнитных сил, возникающих при изменении входного сигнала (например, при превышении уставки), резко изменяли свой выходной сигнал (срабатывали) [62]. Эта особенность и привела к появлению термина «релейная защита» для данного вида автоматики [63].\par\vspace{0.2cm}

% DISS-005-B1474
Достижения в области микроэлектроники и внедрение микропроцессоров значительно усовершенствовали алгоритмы устройств РЗА, что привело к повышению их производительности, чувствительности и селективности [64]. Важно отметить, что современные устройства обеспечивают более простое хранение осциллограмм событий для последующего анализа и точной настройки в сложных сценариях. Однако определенные проблемы сохраняются. Современные энергосистемы функционируют в условиях значительного разнообразия режимов генерации и потребления, а также в сложных и динамично меняющихся эксплуатационных условиях. Этим условиям не всегда могут соответствовать фиксированные настройки РЗА, что снижает эффективность их работы [65, 66, 67].\par\vspace{0.2cm}

% DISS-005-B1475
Перспективным решением этих проблем является использование адаптивных уставок, которые автоматически подстраиваются под текущие условия эксплуатации [68, 69, 70, 71, 72, 73, 74]. Такой подход уменьшает количество параметров, требующих предварительного расчета, и повышает стабильность и производительность устройств РЗА. Однако в некоторых случаях предварительная настройка гибких параметров все же необходима для обеспечения корректной работы системы. Разработка фиксированной логики, учитывающей все возможные режимы и изменения нагрузки, особенно с учетом современных тенденций, таких как расширение сетей и интеграция распределенной генерации, остается сложной задачей [75, 76, 77, 78, 79]. В некоторых случаях несоответствия в поведении устройств РЗА проектным требованиям становятся очевидны только в ходе эксплуатационного анализа. Поэтому для всестороннего тестирования, верификации и дальнейшего совершенствования алгоритмов, включая адаптивные уставки, в частности для устройств быстродействующего автоматического ввода резерва (БАВР), требуются большие объемы репрезентативных данных, таких как реальные осциллограммы и накопленный опыт эксплуатации.\par\vspace{0.2cm}

% DISS-005-B1476
Другим многообещающим направлением является использование методов машинного обучения (ML), обладающих высокой способностью к обобщению на новых данных. ML-модели могут помочь в выявлении релевантных событий или служить независимыми компонентами в адаптивных настройках, или даже функционируя как полноценные элементы РЗА [80, 81]. Последние достижения в области искусственного интеллекта (ИИ), особенно в нейронных сетях, позволяют создавать алгоритмы, способные напрямую обрабатывать сырые данные. Это устраняет необходимость в сегментации данных, извлечении признаков и их отборе, тем самым обеспечивая обнаружение нарушений и их причин в режиме реального времени. Однако обучение таких моделей также требует большого объема данных [82].\par\vspace{0.2cm}

% DISS-005-B1477
Таким образом, создание обширных баз данных осциллограмм для различных процессов является необходимой задачей. Такие базы данных позволяют исследовать новые алгоритмы РЗА, основанные на традиционных принципах, и тестировать их функциональность. Кроме того, они способствуют разработке ML-моделей для решения многочисленных задач, включая модернизацию алгоритмов БАВР.\par\vspace{0.2cm}

% DISS-005-B1478
Многие исследования по применению ML в релейной защите [83, 84] и анализе осциллограмм используют синтетические данные, сгенерированные на основе математических моделей. Такой подход обеспечивает достаточно большие обучающие наборы данных и имеет ряд преимуществ:\par\vspace{0.2cm}

\begin{enumerate}[label=\arabic*),itemsep=0pt,topsep=2pt]
  % DISS-005-B1479
  \item Позволяет генерировать данные в уже унифицированном и нормализованном виде, что упрощает их дальнейшее использование;
  % DISS-005-B1480
  \item Облегчает создание многочисленных вариаций интересующих событий, заданных автором (например, междуфазные замыкания, переходное сопротивление, местоположение повреждения и т.д.);
  % DISS-005-B1481
  \item Обеспечивает контролируемые эксперименты для изучения влияния конкретных факторов;
  % DISS-005-B1482
  \item Поддерживает быструю и простую настройку параметров и масштабирование объемов исследований для конкретных аспектов без дополнительных финансовых затрат;
  % DISS-005-B1483
  \item Позволяет создавать сбалансированные наборы данных.
\end{enumerate}\par\vspace{0.2cm}

% DISS-005-B1484
Следовательно, синтетические данные ценны для изучения применимости методов ML [85] и обучения начальных версий моделей, корректно функционирующих в идеальных условиях. Однако практическое применение требует валидации на других данных различными исследователями, поскольку нейронные сети могут быть очень чувствительны к шуму [86, 87]. Действительно, использование более детализированных математических моделей, например, моделей измерительных трансформаторов, учитывающих насыщение сердечника, позволяет приблизить синтетические данные к реальным. Но, несмотря на преимущества синтетических данных, все еще существуют пробелы. В частности, синтетические данные не могут одновременно учитывать широкий спектр факторов, включая:\par\vspace{0.2cm}

\begin{enumerate}[label=\arabic*),itemsep=0pt,topsep=2pt]
  % DISS-005-B1485
  \item Погрешности и характеристики устройств регистрации событий, а также возможные неисправности в них;
  % DISS-005-B1486
  \item Все нелинейные свойства элементов энергосистемы;
  % DISS-005-B1487
  \item Случайный характер аварийных процессов (например, удары молний, схлестывание проводов, дуговые перемежающиеся замыкания);
  % DISS-005-B1488
  \item Различные типы нагрузок, каждая из которых имеет свои уникальные проблемы моделирования (например, асинхронные и синхронные двигатели, печи, частотные приводы).
\end{enumerate}\par\vspace{0.2cm}

% DISS-005-B1489
Поэтому датасеты, основанные на реальных осциллограммах, имеют большое значение. Однако сбор реальных данных для всех возможных случаев практически невозможен, даже с использованием различных методов аугментации. Таким образом, перспективным является сочетание обоих подходов. Реальные данные могут использоваться для финального тестирования предварительно обученных моделей и для дообучения моделей, изначально обученных на синтетических данных, тем самым повышая их эксплуатационную надежность.\par\vspace{0.2cm}

% DISS-005-B1490
Такие ML-модели могут быть внедрены в промышленные устройства, например, в терминалы РЗА, для улучшения их производительности и облегчения задач персонала. Например, алгоритмы классификации осциллограмм могут предоставлять предварительный анализ аварийных событий и передавать отчеты обслуживающему персоналу. Кроме того, набор реальных данных позволяет верифицировать симуляционные модели, используемые для генерации синтетических данных, облегчает анализ процессов в энергосистемах и способствует созданию новых теорий (например, моделей дуговых перемежающихся однофазных замыканий на землю).\par\vspace{0.2cm}

% DISS-005-B1491
Таким образом, создание базы данных реальных осциллограмм является актуальной задачей для продвижения методов ML в электроэнергетике, верификации симуляционных моделей, генерирующих синтетические данные, и улучшения алгоритмов РЗА, основанных на традиционных методах. Однако на сегодняшний день такие датасеты практически отсутствуют в публичном доступе, что ограничивает развитие подходов, основанных на данных, в этой области. Для устранения этого пробела нами был создан обширный и общедоступный датасет реальных осциллограмм.\par\vspace{0.2cm}

% DISS-005-B1492
Ключевые результаты данного раздела включают:\par\vspace{0.2cm}

\begin{itemize}[itemsep=0pt,topsep=2pt]
  % DISS-005-B1493
  \item \textbf{Обширный датасет реальных осциллограмм:} Представлен общедоступный датасет, включающий 50~765\commentnote{DISS-005-C0105}{Алексей Евдаков [2]}{Может изменяться} реальных осциллограмм, которые охватывают широкий спектр аварийных событий, различных эксплуатационных сценариев, а также записи режимов без существенных событий (шумовые записи)
  % DISS-005-B1494
  \item \textbf{Стандартизированный конвейер обработки данных:} Мы предлагаем систематическую методологию сбора, унификации, анонимизации и нормализации осциллограмм.
\end{itemize}\par\vspace{0.2cm}

% DISS-005-B1495
В настоящей главе подробно описываются этапы сбора, предобработки и подготовки датасета для его последующего использования в задачах, связанных с анализом и модернизацией устройств БАВР. Особое внимание уделяется методам унификации, деперсонализации и нормализации данных, а также подходам к их разметке для обучения моделей машинного обучения.\par\vspace{0.2cm}

% DISS-005-B1496
В следующих главах (\textbf{ГЛАВА ХХ}\commentnote{DISS-005-C0106}{Алексей Евдаков [2]}{Вставить ссылку}) будет разобраны варианты применения систем на основе нейронных сетей. Мы разрабатываем автоматизированную систему, использующую нейронные сети для классификации и сортировки осциллограмм. Так же можно будет указать иные анализы данных\par\vspace{0.2cm}

% DISS-005-B1497

% DISS-005-B1498
\section{Методология формирования датасета реальных осциллограмм}\label{sec:dataset_methodology}

% DISS-005-B1499
\subsection{Этапы подготовки датасета}\label{subsec:dataset_stages}

% DISS-005-B1500
Создание комплексного и пригодного для использования датасета реальных осциллограмм является многоэтапным процессом, требующим систематического подхода. Классические этапы подготовки обширной базы данных обычно включают следующие шаги [88, 89, 90, 91], которые были адаптированы и детализированы с учетом специфики поставленной задачи:\par\vspace{0.2cm}

\begin{enumerate}[label=\arabic*),itemsep=0pt,topsep=2pt]
  % DISS-005-B1501
  \item \textbf{Сбор\commentnote{DISS-005-C0107}{Алексей Евдаков [2]}{Есть неплохая идея от нейронки про блок схему отражающий сей порядок и следующие разделы. Она там предложила такой порядок:\newline 1. Сбор данных --- 50~765 файлов из X источников\newline 2. Проверка форматов/целостности\newline 3. Унификация форматов (COMTRADE) и имён\newline 4. Деперсонализация\newline 5. Нормализация\newline 6. Разметка\newline 7. Экспорт в CSV}\commentnote{DISS-005-C0108}{Алексей Евдаков [2]}{Может потом сделаю, пока это как мысль.} данных из различных источников:} Поиск и агрегация осциллограмм с различных энергообъектов, от разных производителей устройств РЗА и из различных эксплуатационных ситуаций.
  % DISS-005-B1502
  \item \textbf{Предварительная очистка данных:} Начальная фильтрация, удаление заведомо некорректных или нечитаемых файлов, проверка целостности записей (например, наличие пары .cfg и .dat для формата COMTRADE).
  % DISS-005-B1503
  \item \textbf{Унификация форматов и стандартизация имен сигналов:} Приведение осциллограмм к единому формату (в нашем случае, к ASCII COMTRADE IEEE C37.111-2013) и разработка единой системы именования для аналоговых и дискретных сигналов, что критически важно для последующей автоматизированной обработки и анализа.
  % DISS-005-B1504
  \item \textbf{Деперсонализация (анонимизация) данных:} Удаление любой конфиденциальной информации, которая могла бы идентифицировать конкретный энергообъект, время события или производителя работ, для обеспечения возможности публичного использования датасета. Отметим, что для внутреннего использования (например, для анализа работы оборудования на конкретном объекте) исходная информация может быть сохранена. Однако для создания публичного датасета, как описано в данной работе, все данные проходят процедуру анонимизации.
  % DISS-005-B1505
  \item \textbf{Нормализация сигналов:} Приведение амплитуд сигналов к сопоставимым величинам (например, к относительным единицам или стандартизированным вторичным значениям) для корректного сравнения и использования в алгоритмах, чувствительных к масштабу данных. Хотя данные от измерительных трансформаторов часто приводятся к стандартным вторичным значениям, в файлах осциллограмм не всегда явно указаны коэффициенты трансформации или тип величин (первичные/вторичные). В ряде современных устройств, особенно использующих цифровые измерительные датчики, понятие классических вторичных значений может отсутствовать. Это создает специфические трудности при работе с реальными осциллограммами и требует разработки подходов к их корректной нормализации.
  % DISS-005-B1506
  \item \textbf{Разметка (аннотирование) осциллограмм:} Классификация событий, зафиксированных в осциллограммах, определение их временных границ и присвоение меток. Этот этап может быть как ручным, так и частично или полностью автоматизированным. Этот этап может быть как ручным, так и частично или полностью автоматизированным. В рамках данной работы для части датасета применялась ручная разметка, а также исследовались подходы к автоматизированной разметке.
  % DISS-005-B1507
  \item \textbf{Документирование и описание датасета:} Создание подробного описания структуры датасета, форматов данных, методологии сбора и обработки, а также характеристик включенных осциллограмм. Что будет подробно разобрано далее
  % DISS-005-B1508
  \item \textbf{Публикация и распространение (для публичного датасета):} Обеспечение доступа к датасету для научного сообщества. Информация о доступе представлена по ссылкам [ХХ\commentnote{DISS-005-C0109}{Алексей Евдаков [2]}{Появится в будущем}] и в публикациях [ХХ\commentnote{DISS-005-C0110}{Алексей Евдаков [2]}{Тоже появится в будущем}]
  % DISS-005-B1509
  \item \textbf{Обновление и поддержка:} Периодическое пополнение датасета новыми данными и актуализация документации.
\end{enumerate}\par\vspace{0.2cm}

% DISS-005-B1510
В рамках настоящей работы основное внимание уделяется этапам 1-7. Вопросы публикации и долгосрочной поддержки выходят за рамки данной диссертации, однако их важность для создания ценного научного ресурса несомненна. Последующие разделы данной главы детально описывают реализацию этих этапов применительно к созданию нашего датасета осциллограмм\par\vspace{0.2cm}

% DISS-005-B1511
\subsection{Проблемы и вызовы при работе с реальными осциллограммами\commentnote{DISS-005-C0111}{Алексей Евдаков [2]}{В теории – МЭК 61850 направлен на решение этого. Вот только это идеи о будущем\dots\ Возможно стоит как-то про это сказать.}}\label{subsec:dataset_challenges}

% DISS-005-B1512
Работа с реальными осциллограммами, в отличие от симулированных данных, сопряжена с рядом специфических проблем и вызовов, требующих тщательного учета и разработки специализированных методов их преодоления:\par\vspace{0.2cm}

\begin{itemize}[itemsep=0pt,topsep=2pt]
  % DISS-005-B1513
  \item \textbf{Разнообразие форматов файлов осциллограмм:} Устройства РЗА от различных производителей могут записывать осциллограммы во множестве проприетарных или стандартизированных форматов (например, COMTRADE различных версий~-- «.cfg+.dat», «.cff»; форматы НПП ЭКРА~-- «.dfr»; НПП Бреслер~-- «.brs»; НЕВА~-- «.os», «.os1»; ПАРМА~-- «.do», «.to»; регистраторы ЧЭАЗ~-- «.osc»; ПРОСОФТ~-- «.sg2»; регистраторы переходных процессов~-- «.bb» и др.). Это требует разработки или использования универсальных инструментов для чтения и конвертации данных.
  % DISS-005-B1514
  \item \textbf{Необходимость унификации параметров:}
  \begin{itemize}[itemsep=0pt,topsep=2pt]
    % DISS-005-B1515
    \item Список и названия аналоговых и дискретных сигналов: Имена сигналов часто назначаются производителем терминала или инженером при наладке и могут сильно варьироваться даже для однотипных измерений на разных объектах или устройствах. Отсутствие единого стандарта именования затрудняет автоматизированный анализ и сопоставление данных.
    % DISS-005-B1516
    \item Частота дискретизации и частота сети: Осциллограммы могут быть записаны с различной частотой дискретизации (например, 1000, 1200, 1600, 2400, 6400~Гц и т.д.) и для сетей с разной номинальной частотой (50~Гц, 60~Гц). Эти параметры напрямую влияют на количество точек на период и требуют учета при спектральном анализе или применении алгоритмов, основанных на временных окнах.
    % DISS-005-B1517
    \item Стандартные значения амплитуд (коэффициенты трансформации): Информация о коэффициентах трансформации измерительных трансформаторов тока и напряжения (или других датчиков) не всегда явно присутствует или легко извлекаема из файла осциллограммы, что усложняет процесс нормализации и приведения к первичным или стандартизированным вторичным величинам.
  \end{itemize}
  % DISS-005-B1518
  \item \textbf{Большой объем данных и трудоемкость ручной обработки:} Современные системы РЗА генерируют значительные объемы осциллограмм. Ручная обработка, анализ и разметка такого массива данных практически невозможны и требуют разработки автоматизированных подходов.
  % DISS-005-B1519
  \item \textbf{Обеспечение конфиденциальности и необходимость деперсонализации:} Реальные осциллограммы часто содержат информацию, являющуюся коммерческой тайной или позволяющую идентифицировать конкретный энергообъект (названия подстанций, фидеров, время события, данные о производителе оборудования). Перед любой публикацией или передачей третьим лицам такие данные должны быть тщательно деперсонализированы (анонимизированы) [92, 93, 94].
  % DISS-005-B1520
  \item \textbf{«Зашумленность» данных и наличие неинформативных осциллограмм:} Значительная часть регистрируемых осциллограмм может не содержать полезной информации о значимых событиях (так называемые «пустые» осциллограммы, содержащие только фоновый шум или незначительные флуктуации). Хотя такие данные могут быть полезны для обучения моделей отличать нормальный режим от аномального, их большое количество усложняет поиск и анализ действительно информативных записей. Кроме того, реальные сигналы всегда содержат различные виды помех и наводок.
  % DISS-005-B1521
  \item \textbf{Неполнота или противоречивость данных:} В некоторых осциллограммах могут отсутствовать важные сигналы, или информация в конфигурационном файле (.cfg) может не полностью соответствовать данным в файле значений (.dat). Также встречаются осциллограммы с поврежденной структурой или некорректной кодировкой.
  % DISS-005-B1522
  \item \textbf{Сложность определения «истинной» картины события:} Одна и та же авария может быть зафиксирована несколькими устройствами РЗА с разных точек энергосистемы, и их «взгляд» на событие может несколько отличаться из-за погрешностей измерений, разных алгоритмов фиксации и т.д.
\end{itemize}\par\vspace{0.2cm}

% DISS-005-B1523
Успешное решение этих проблем является необходимым условием для создания качественного и репрезентативного датасета реальных осциллограмм, пригодного для решения широкого круга задач в области релейной защиты, автоматики и машинного обучения.\par\vspace{0.2cm}

% DISS-005-B1524
